\section{Hipótesis}

¿Cómo puede una aplicación móvil de escaneo de código de barra contribuir a reducir 
la asimetría de información en la compra de productos electrónicos y fomentar un 
consumo más informado y sostenible? 

Mediante el desarrollo de una aplicación que centralice la información actualmente 
dispersa sobre productos electrónicos —con énfasis en smartphones y laptops— y sus 
fabricantes, se espera generar en el consumidor una reflexión crítica sobre sus propios
hábitos de consumo: que se cuestione cuáles son realmente sus necesidades como persona 
antes de asumirlas como necesidades de consumidor, pues identificar ese origen es el 
primer paso para tomar decisiones conscientes. 

Se espera que la herramienta impulse la adopción de principios de la economía circular, 
dotando al usuario de información suficiente para prevenir el desperdicio de tiempo y dinero. 
Esto implica que pueda evaluar la calidad real de un producto a partir de sus especificaciones 
técnicas y reseñas, compararlo con otras opciones disponibles en el mercado, y valorarlo según 
su comportamiento social y su capacidad de reparación. Este último factor —la reparabilidad— 
puede transitar de ser un aspecto técnico complejo y poco conocido a convertirse en un criterio 
accesible y popular; condición que amplía la diversidad de soporte disponible para el dispositivo 
y fortalece la posición del consumidor frente a la narrativa que las empresas construyen por conveniencia. 

Se espera que la aplicación permita que cualquier persona, con tan solo escanear el código 
de barras de un producto, acceda a información integral sobre el mismo. Esta no se limitará 
a sus especificaciones técnicas, sino que también incluirá datos relevantes sobre la empresa 
fabricante: sus sedes, certificaciones, alianzas comerciales y su compromiso real con el medio 
ambiente. Adicionalmente, la plataforma facilitará la localización de centros de reparación autorizados 
o cercanos al usuario, fomentando la reutilización de dispositivos y extendiéndoles una segunda vida 
útil. Por último, y como elemento vital, integrará reseñas de consumidores del producto: opiniones 
no patrocinadas, reales y conscientes. 
